\section{KARŞILAŞILAN ZORLUKLAR}

Projenin geliştirme sürecinde donanım ve yazılım bileşenlerinin aynı sistemde kararlı çalıştırılması bazı zorluklar oluşturmuştur. Özellikle servo motorlar, analog mesafe sensörü, I2C cihazları, SPI tabanlı SD kart modülü ve web arayüzü birlikte kullanıldığında güç, haberleşme ve zamanlama kaynaklı problemler görülmüştür.

\subsection{Güç Besleme ve Servo Kaynaklı Reset Problemleri}
İlk testlerde sistemin zaman zaman durduğu ve Arduino'nun yeniden başladığı gözlemlenmiştir. Bunun temel nedeni, servo motorların hareket başlangıcında yüksek akım çekmesidir. Arduino 5V hattı bu akımı karşılayamadığında besleme gerilimi düşmekte ve sistem reset atabilmektedir.

Bu nedenle servo motorların Arduino 5V pininden değil, harici bir güç kaynağından beslenmesi daha uygun görülmüştür. Harici besleme kullanıldığında Arduino GND hattı ile servo beslemesinin GND hattı ortak bağlanmalıdır. Aksi durumda servo sinyali doğru referans alamaz ve motor hareketleri kararsızlaşır.

\subsection{Sensör Ölçüm Gürültüsü ve Hedef Kararı}
Sharp GP2Y0A02YK0F analog çıkışlı bir sensördür. Bu nedenle ölçümler ışık koşulu, hedef yüzeyi, hedef açısı ve besleme kararlılığından etkilenebilir. Sensörden gelen her değerin doğrudan hedef olarak kabul edilmesi yanlış sonuçlara yol açabilir.

Bu problemi azaltmak için ölçüm değerleri mesafeye dönüştürülmüş, geçersiz aralıklar filtrelenmiş ve belirlenen eşik değerinin altındaki ölçümler hedef olarak kabul edilmiştir.

\subsection{I2C ve SPI Cihazlarının Birlikte Kullanımı}
Projede LCD ekran ve MPU6050 I2C hattını, SD kart modülü ise SPI hattını kullanmaktadır. Birden fazla haberleşme protokolünün aynı Arduino üzerinde çalışması bağlantı ve zamanlama açısından dikkat gerektirmiştir.

LCD adresinin 0x3F, MPU6050 adresinin 0x68 olması I2C adres çakışmasını önlemiştir. SD kart için CS pini D4 olarak belirlenmiş, MOSI, MISO ve SCK pinleri Arduino Uno'nun donanımsal SPI pinlerine bağlanmıştır. exFAT desteği gerektiği için standart SD kütüphanesi yerine SdFat kütüphanesi tercih edilmiştir.

\subsection{Seri Haberleşme ve Web Arayüzü Problemleri}
Arduino'dan web arayüzüne JSON formatında veri gönderilmektedir. Seri haberleşmede bazı paketler eksik veya parçalı gelebilmektedir. Bu durumda web arayüzünde “geçersiz veri” kayıtları oluşmuştur.

Çözüm olarak web tarafında satır sonuna göre paket biriktirme yapılmış, yalnızca tamamlanmış JSON satırları işlenmiştir. Hatalı paketler sistemin çalışmasını durdurmadan yok sayılmıştır. Ayrıca olay kaydı alanı sınırlandırılarak nesne algılandığında sayfanın aşağı kayması engellenmiştir.

\section{SONUÇLAR VE ÖNERİLER}

Bu projede, düşük maliyetli donanımlar kullanılarak pan-tilt kızılötesi radar prototipi başarıyla gerçekleştirilmiştir. Sistem; mesafe algılama, servo hareketi, LCD gösterimi, MPU6050 takibi, SD kart veri kaydı ve web arayüzünde gerçek zamanlı radar görüntüleme işlevlerini bir arada sunmaktadır.

Elde edilen sonuçlar, Arduino Uno ile birden fazla çevresel birimin aynı anda yönetilebileceğini göstermiştir. Ancak kararlı çalışma için güç beslemesi, ortak topraklama, kablo bağlantıları, sensör kalibrasyonu ve seri veri paketleme yapısı kritik öneme sahiptir. Özellikle servo motorların ani akım çekimi sistem resetlerinin temel nedeni olarak değerlendirilmiştir.

Web arayüzü, sistem verilerinin daha anlaşılır biçimde izlenmesini sağlamıştır. Pan açısı, tilt açısı, mesafe ve hedef durumu aynı ekranda gösterilerek test süreci kolaylaştırılmıştır. SD kart desteği ise ölçümlerin sonradan incelenebilmesi için yararlı bir altyapı sunmuştur.

Gelecek çalışmalar için öneriler aşağıda verilmiştir:

\begin{itemize}
    \item Sharp sensör için daha ayrıntılı kalibrasyon yapılabilir.
    \item Analog ölçüm gürültüsü için hareketli ortalama veya medyan filtre eklenebilir.
    \item Servo motorlar için daha kararlı ve regüle edilmiş ayrı bir güç kaynağı kullanılabilir.
    \item Pan-tilt mekanizması daha rijit bir gövde ile geliştirilebilir.
    \item Kızılötesi sensör yerine LiDAR, ToF veya ultrasonik sensörler denenebilir.
    \item MPU6050 verileri titreşim telafisi için kontrol algoritmasına dahil edilebilir.
    \item SD karta kaydedilen CSV verilerinden performans analizi yapılabilir.
    \item Web arayüzüne geçmiş taramaları görüntüleme ve veri indirme özellikleri eklenebilir.
    \item Arduino Uno yerine ESP32 gibi daha güçlü ve kablosuz haberleşme destekli kartlar kullanılabilir.
\end{itemize}

Sonuç olarak geliştirilen sistem, gömülü sistem tasarımında sensör okuma, eyleyici kontrolü, haberleşme, güç yönetimi ve gerçek zamanlı arayüz geliştirme konularını bir araya getiren uygulanabilir bir prototip olmuştur.
